\section{Cable paths and the joint limit}\label{sec:paths}

Proposition~\ref{prop:path-limit} identifies the walk in conformal
coordinates and with exponential holding times.  We now place the original
cable graphs and the limiting metric sphere in one compact space, replace
the exponential clock by the deterministic clock, and restore the
randomness of the environment.  Until the final proof, the realization
supplied by Proposition~\ref{prop:inputs} remains fixed.

\subsection{A common metric space}

The following construction uses only the vertex-distance comparison.  In
particular, it does not require an extension of the conformal coordinate
map to the interiors of cable edges.

\Needspace{8\baselineskip}
\begin{lemma}[A compact ambient space]\label{lem:ambient-space}
Let $C_n$ be finite connected cable graphs, with vertex sets $V_n$ and
all edge lengths equal to $a_n>0$, where $a_n\to0$.  Write $d_{C_n}$ for
their path metrics.  Let $(S,d)$ be a compact metric space, and suppose
that maps $z_n:V_n\to S$ satisfy
\begin{align*}
 \delta_n&:=\sup_{u,v\in V_n}
  \abs{d_{C_n}(u,v)-d(z_n(u),z_n(v))}\longrightarrow0,\\
 \varepsilon_n&:=d_{\mathrm H}^{S}(z_n(V_n),S)\longrightarrow0.
\end{align*}
There is a compact metric space $(Z,d_Z)$ containing isometric copies of
$S$ and every $C_n$, such that, with $c_n=\delta_n/2+a_n$,
\begin{align}
 \sup_{v\in V_n}d_Z(v,z_n(v))&\le c_n,
 \label{eq:vertex-bridges}\\
 d_{\mathrm H}^{Z}(C_n,S)&\le
  \max\{a_n+c_n,\varepsilon_n+c_n\}.
 \label{eq:cable-hausdorff}
\end{align}
If $\nu_n$ is a probability measure on $V_n$ and
$(z_n)_*\nu_n\Rightarrow\nu$ on $S$, then
$d_{\mathrm P}^{Z}(\nu_n,\nu)\to0$.
\end{lemma}

\begin{proof}
Start with the disjoint union of $S$ and the $C_n$.  Retain the metric
within each piece, and allow a crossing from $v\in V_n$ to $z_n(v)\in S$
at cost $c_n$.  Define distance as the infimum of the costs of finite
chains made of such crossings and movements within individual pieces.

This construction does not shorten any of the original metrics.  Indeed,
an excursion from $S$ through $C_n$, entering at $z_n(u)$ and returning
at $z_n(v)$, costs at least
\[
 2c_n+d_{C_n}(u,v)\ge d(z_n(u),z_n(v)).
\]
Replacing all such excursions by movements in $S$ proves the assertion
for pairs of points of $S$.  For points of a fixed $C_n$, first remove
excursions through the other cable graphs in this way.  Each remaining
excursion out of $C_n$, from a vertex $u$ to a vertex $v$, costs at least
\[
 2c_n+d(z_n(u),z_n(v))\ge d_{C_n}(u,v),
\]
so it too can be removed.  This also proves the assertion for points in
edge interiors.  Distinct pieces remain separated: every chain leaving
$C_n$ pays at least the positive cost $c_n$.  Thus the chain construction
defines a metric, not merely a pseudometric.

Every point of $C_n$ is at distance at most $a_n$ from a vertex, even
when its edge is a loop.  It is consequently at distance at most
$a_n+c_n$ from $S$.  Conversely, every point of $S$ is at distance at
most $\varepsilon_n+c_n$ from $C_n$.  These observations prove
\eqref{eq:vertex-bridges}--\eqref{eq:cable-hausdorff} on the union.  The
union is totally bounded: its tail lies arbitrarily close to the compact
space $S$, and each of the finitely many remaining cable graphs is
compact.  Its completion is the required compact space $Z$.

The coupling $v\mapsto(v,z_n(v))$ of $\nu_n$ and $(z_n)_*\nu_n$ has
all its mass at pairs at distance at most $c_n$, whence
\[
 d_{\mathrm P}^{Z}(\nu_n,\nu)
 \le c_n+d_{\mathrm P}^{Z}((z_n)_*\nu_n,\nu)\longrightarrow0.
\]
For the last convergence, weak convergence on the compact subspace $S$
is also weak convergence in $Z$.
\end{proof}

Apply Lemma~\ref{lem:ambient-space} to
$C_n=(\overline M_n,a_nd_n)$ and $(S,d)=(\Sph,D_h)$, with
$\nu_n=\mu_n$.  The hypotheses are precisely the metric and measure
conclusions of Proposition~\ref{prop:inputs}; vertex distances in a
cable graph equal the corresponding graph distances.  We henceforth
regard these spaces as subsets of the resulting $Z$.

Because $\rho$ and $D_h$ induce the same topology on the compact sphere,
the identity map between the two metric spheres is uniformly continuous
in both directions.  Applying the same time changes in the definition
of the Skorokhod $J_1$ topology therefore transfers
Proposition~\ref{prop:path-limit} from $\rho$ to $D_h$.  Moreover,
\eqref{eq:vertex-bridges} gives the deterministic bound
\[
 \sup_{0\le t\le T}d_Z(Y^n_t,\xi^n_t)\le c_n.
\]
It follows that, for every $T<\infty$,
\begin{equation}\label{eq:ambient-path-limit}
 Y^n\ \Longrightarrow\ B^h
 \quad\text{in }D([0,T],Z)\text{ with the }J_1\text{ topology}.
\end{equation}

We shall use one consequence of the continuity of $B^h$.  For a path
$y:[0,T]\to Z$, set
\[
 \omega_T(y,s)=
 \sup_{\substack{t,t'\in[0,T]\\\abs{t-t'}\le s}}
    d_Z(y_t,y_{t'}).
\]
Then, for every $\eta>0$,
\begin{equation}\label{eq:continuous-limit-modulus}
 \lim_{s\downarrow0}\limsup_n
 \PP\bigl[\omega_T(Y^n,s)>\eta\bigr]=0.
\end{equation}
To see this directly, use the Skorokhod representation theorem
\cite[Theorem~6.7]{Billingsley1999} on the Polish space $D([0,T],Z)$ in
\eqref{eq:ambient-path-limit}.  Convergence in $J_1$ to a continuous
path is uniform: composing with the convergent time changes changes that
limiting path uniformly little; see also
\cite[Section~12, p.~124]{Billingsley1999}.  On this representation,
\[
 \omega_T(Y^n,s)\le
 2\sup_{t\le T}d_Z(Y^n_t,B^h_t)+\omega_T(B^h,s).
\]
Uniform continuity of each limiting path proves
\eqref{eq:continuous-limit-modulus}.  The resulting probability statement
depends only on the original path laws, not on the chosen representation.

\subsection{Removing the exponential clock}

The clock comparison must count every attempted step, including a loop
traversal that does not change the vertex of the continuous-time chain.

\begin{lemma}[The attempt clock]\label{lem:clock}
Fix a realization supplied by Proposition~\ref{prop:inputs}, for which
\eqref{eq:ambient-path-limit} holds.  Let $X^n$ be the stationary discrete
walk on $M_n$, including its selected half-edges, and let $\overline X^n$
be its cable interpolation.  Independently of the walk, let
$(e_{n,k})_{k\ge1}$ be
independent exponential random variables of mean $1/n$, and put
\[
 S_{n,0}=0,\qquad S_{n,k}=\sum_{j=1}^k e_{n,j},\qquad
 N_n(t)=\max\{k:S_{n,k}\le t\}.
\]
Then $Y^n_t=X^n_{N_n(t)}$ has generator $-G_n$ and
$Y^n_{S_{n,k}}=X^n_k$.  In the ambient space of
Lemma~\ref{lem:ambient-space}, for every $T<\infty$,
\begin{equation}\label{eq:clock-comparison}
 \sup_{0\le t\le T}
 d_Z(\overline X^n_{nt},Y^n_t)\longrightarrow0
 \quad\text{in probability}.
\end{equation}
\end{lemma}

\begin{proof}
Attempts occur at rate $n$ and apply the transition kernel $P_n$,
so the generator is $n(P_n-I)=-G_n$.  This remains true when an
attempt traverses a loop.  The identity at $S_{n,k}$ follows from the
right-continuous definition of $N_n$.

Fix $T$ and write
\[
 \Delta_{n,T}=
 \max_{0\le k\le\lfloor nT\rfloor}\abs{S_{n,k}-k/n}.
\]
The centered partial sums form a martingale.  Since
$\operatorname{Var}(e_{n,k})=n^{-2}$, the $L^2$ maximal inequality gives
\begin{equation}\label{eq:clock-deviation}
 \EE[\Delta_{n,T}^2]
 \le \frac{4\lfloor nT\rfloor}{n^2}\longrightarrow0.
\end{equation}
For $k=\lfloor nt\rfloor$, the point $\overline X^n_{nt}$ lies on the
cable edge traversed from $X^n_k$, at distance at most $a_n$ from that
vertex.  Thus
\[
 d_Z(\overline X^n_{nt},Y^n_t)
 \le a_n+d_Z(Y^n_{S_{n,k}},Y^n_t),
 \qquad
 \abs{S_{n,k}-t}\le\Delta_{n,T}+1/n.
\]
The spatial bound holds for loop edges as well as nonloop edges.
On the event $\Delta_{n,T}+1/n\le s\le1$, all the times appearing here
lie in $[0,T+1]$.  Consequently, for $a_n<\eta/2$,
\[
 \PP\left[\sup_{t\le T}
 d_Z(\overline X^n_{nt},Y^n_t)>\eta\right]
 \le \PP[\Delta_{n,T}+1/n>s]
   +\PP[\omega_{T+1}(Y^n,s)>\eta/2].
\]
First take the upper limit in $n$, using
\eqref{eq:clock-deviation}, and then let $s\downarrow0$, using
\eqref{eq:continuous-limit-modulus} on $[0,T+1]$.  This proves
\eqref{eq:clock-comparison}.
\end{proof}

Combining Lemma~\ref{lem:clock} and \eqref{eq:ambient-path-limit} gives
convergence of the cable paths to $B^h$ in $D([0,T],Z)$.  Both the cable
paths and the limiting paths are continuous.  The continuous path space
is closed in the $J_1$ topology: a convergent sequence of continuous
paths, after the defining time changes, converges uniformly and hence
has a continuous limit.  Its induced topology is the uniform topology,
by the same time-change argument used above.  Thus in fact
\begin{equation}\label{eq:cable-finite-horizon}
 (\overline X^n_{nt})_{0\le t\le T}
 \ \Longrightarrow\ (B^h_t)_{0\le t\le T}
 \quad\text{in }C([0,T],Z).
\end{equation}

This convergence holds on every integer time horizon, and therefore also
in $C([0,\infty),Z)$ with the topology of uniform convergence on compact
time intervals.  Here are the compactness details.  For each integer
$m\ge1$ and $\epsilon>0$, finite-horizon tightness supplies a compact
set $K_m\subset C([0,m],Z)$ such that the probability of leaving $K_m$
is at most $\epsilon2^{-m}$, uniformly in $n$.  The set of paths whose
restriction belongs to $K_m$ for every $m$ is compact in
$C([0,\infty),Z)$, by a diagonal subsequence argument, and has
probability at least $1-\epsilon$.  Every subsequential limit has, by
the continuous restriction maps, the law of $B^h$ on every $[0,m]$.
These restrictions determine the entire path law.  We have proved
\begin{equation}\label{eq:cable-local-limit}
 (\overline X^n_{nt})_{t\ge0}
 \ \Longrightarrow\ (B^h_t)_{t\ge0}
 \quad\text{in }C([0,\infty),Z).
\end{equation}

\subsection{Restoring the random environment}

We have obtained the conditional path limit for every fixed realization
in the probability-one event of Proposition~\ref{prop:inputs}.  We now
combine it with convergence of the underlying measured spaces.

\begin{proof}[Proof of Theorem~\ref{thm:main}]\label{proofmain}
Start from an arbitrary subsequence and take the coupled further
subsequence of Proposition~\ref{prop:inputs}.  Fix a realization in
its probability-one event and a finite integer $r\ge1$.
Conditional on this realization, the $r$ cable paths are independent
with the common law in \eqref{eq:cable-local-limit}.  Their joint law
therefore converges to that of $r$ independent copies of $B^h$ in
$C([0,\infty),Z)^r$.  For example, a Skorokhod representation of
\eqref{eq:cable-local-limit}, repeated on $r$ independent probability
spaces, couples the product laws with almost sure convergence of every
coordinate.

Lemma~\ref{lem:ambient-space} also gives Hausdorff convergence of the
spaces and Prokhorov convergence of their measures in the same $Z$.
Under the preceding product coupling, the curves converge uniformly on
every compact time interval.  By the definition of the curve-decorated
Gromov--Hausdorff--Prokhorov topology, the decorated spaces consequently
converge almost surely under that coupling.  In particular, for any
bounded continuous function $F$ on the space of such decorated objects,
the expectation of $F$ over the $r$ walks on this fixed environment
converges to its expectation over $r$ independent stationary Liouville
Brownian motions on the fixed limiting sphere.

These expectations are the conditional expectations obtained by sampling
walks with fresh randomness after the environment coupling.  The walk
kernel on a finite graph depends only on the map, not on the auxiliary
coordinates, marks, or tours used in that coupling; the limiting kernel
is the diffusion law specified by the quantum sphere.  We have proved
pointwise convergence of these conditional expectations on the
probability-one event.  The compact space $Z$, the path couplings, and
the recovery functions used earlier are witnesses for this pointwise
statement.  They need not be chosen measurably as the environment varies.
The conditional expectations from these sampling kernels are measurable,
and bounded convergence now gives convergence of the unconditional
expectations of $F$ along the coupled further subsequence.

Thus every subsequence admits a further subsequence with the asserted
joint limit.  This proves convergence along the full sequence, for
each fixed $r$.  Lemma~\ref{lem:clock} used the deterministic
acceleration $b_n=n$, which has exactly the time normalization in the
theorem.
\end{proof}
